\chapter{概念设计}
% 各模块概念设计按照：机械、电控、视觉三部分组织，便于后续详细设计和验证逻辑衔接。

\section{机械概念设计}
在比赛中，考虑到比赛中我们希望XX有较高的机动性能甚至实现飞坡的要求， 则要求XX机器人的底盘需要满足\tabref{tab:XX底盘功能设计}中的指标。

\subsection{底盘功能设计}

\begin{longtable}[htbp!]{ m{4cm}<{\centering}  m{9cm}<{\centering} }
\hline
种类  &  指标 \\
\hline
\endfirsthead
\hline
种类  &  指标 \\
\hline
\endhead
\hline
\multicolumn{2}{r}{接下页}
\endfoot
\caption{XX底盘功能设计}
\label{tab:XX底盘功能设计}
\endlastfoot
%   填这里
 \hline
 轮组  & 结构强度高，便于拆装与维护 \\  
 \hline 
 舵结构 & 结构强度高，转动惯量小，精准复位 \\  
 \hline
 悬挂系统  & 硬度可调，可靠性高，易于维护  \\  
 \hline
 底盘车架  & 强度高，重量轻  \\
 \hline
 轮组系统整体  & 结构简洁，内部摩擦损耗小   \\
 \hline
 底盘整体  & 尺寸紧凑，重心低，强度高，轻量化，对称分布  \\
 \hline
\end{longtable}